%%% -*-LaTeX-*-
% arara: make: { }
%%% A four-pass solution without make, and without dependency checks,
%%% might be done like this:
%%: % arara: clean:     { extensions: [ aux, bbl, idx, ind, lof, log, lot, out, stripped, toc ] }
%%:
%%: % arara: pdflatex:  { }
%%: % arara: bibtex:    { }
%%: % arara: makeindex: { }
%%:
%%: % arara: pdflatex:  { }
%%: % arara: bibtex:    { }
%%: % arara: makeindex: { }
%%:
%%: % arara: pdflatex:  { }
%%: % arara: bibtex:    { }
%%: % arara: makeindex: { }
%%:
%%: % arara: pdflatex:  { }
%%% ====================================================================
%%% This is a sample top-level LaTeX file for typesetting a dissertation
%%% or thesis at the University of Utah in a new design in
%%% 2018 that should permit considerably more flexibility and power in
%%% thesis formatting, while supplying the necessary university-specific
%%% front matter that is a part of any dissertation or thesis.  The
%%% guiding principle of the design is that most standard LaTeX packages
%%% should be usable as extensions to the standard LaTeX book class.
%%% In particular, color, hypertext links in PDF output, indexing,
%%% new float types, optional (starred) arguments in LaTeX macros, PDF
%%% document inclusions (for published journal articles), and running
%%% page headers, are all possible, but optional.
%%%
%%% All style files for typesetting University of Utah dissertations
%%% and theses are prefixed with "uu", and users should NEVER introduce
%%% additional personal styles with that prefix.
%%%
%%% Most students find it convenient to start with a COPY of this file
%%% as a template, and then alter that copy to match their needs.
%%% Please RENAME the sample-thesis-*.* files using a suitable prefix
%%% for YOUR thesis, and adjust the Makefile accordingly.
%%%
%%% There is an associated Unix Makefile that can be similarly
%%% customized, and then the only command ever needed to typeset the
%%% complete thesis is the single word "make".  Of course, during
%%% writing and typesetting, not all of the steps are needed, so
%%% often, one can just name a convenience target such as "make
%%% dvi-pass" or "make pdf-pass" to do just a single pass of LaTeX and
%%% BibTeX.
%%%
%%% There should be no, or very few, macro definitions in this file;
%%% any needed belong in a private style file, perhaps called
%%% mythesis.sty, and input below after all other packages.  The bulk
%%% of this file should just be command invocations, and any arguments
%%% that they need.
%%%
%%% We exploit the fact that TeX ignores spaces after command names to
%%% line up arguments for better readability.
%%%
%%% Each chapter should be a separate complete file, so that you can
%%% insert a command like \includeonly{chap_intro} before the first
%%% \include{chap_xxx} command to avoid typesetting all but the
%%% chapter that you are currently working on, to save time.
%%%
%%% Remember that occupants of job positions change jobs from time to
%%% time: YOU are responsible for ensuring correct names of all humans
%%% mentioned in this file!
%%%
%%% In the rest of this top-level thesis file, comments that contain
%%% the word OPTIONAL hide commands that you might wish to use; just
%%% remove the "%%: OPTIONAL: " comment prefix to enable them.
%%%
%%% [25-Jan-2018]
%%% ====================================================================

\documentclass[
    11pt                                % One of: <empty>, or 11pt, or 12pt
]{book}  % everything after this extends the standard LaTeX book class

\overfullrule = 0pt                     % do not show overfull boxes
\overfullrule = 100pt                   % make overfull boxes really visible!

\usepackage {uuthesis-base}             % MANDATORY package for a minimal thesis

%%% ====================================================================
%%% Choose an alternate font family for the document if the TeX default
%%% of Computer Modern is not wanted.  There are several packages that
%%% provide complete TeX math fonts in attractive type designs.  Pick
%%% JUST ONE of them, or leave them all commented out for the default.

%%: OPTIONAL:  \usepackage {bookman}
%%: OPTIONAL:  \usepackage {caladea}
%%: OPTIONAL:  \usepackage {charter}
%%: OPTIONAL:  \usepackage {cochineal}
%%: OPTIONAL:  \usepackage {concrete}
%%: OPTIONAL:  \usepackage {crimson}
%%: OPTIONAL:  \usepackage {fourier}
%%: OPTIONAL:  \usepackage {gfsartemisia}
%%: OPTIONAL:  \usepackage {heuristica}
%%: OPTIONAL:  \usepackage {kerkis}
%%: OPTIONAL:  \usepackage {librebaskerville}
%%: OPTIONAL:  \usepackage {LibreBodoni}
%%: OPTIONAL:  \usepackage {librecaslon}
%%: OPTIONAL:  \usepackage {mathpazo}
%%: OPTIONAL:  \usepackage {merriweather}
%%: OPTIONAL:  \usepackage {tgadventor}
%%: OPTIONAL:  \usepackage {tgbonum}
%%: OPTIONAL:  \usepackage {tgcursor}
%%: OPTIONAL:  \usepackage {tgheros}
%%: OPTIONAL:  \usepackage {tgpagella}
%%: OPTIONAL:  \usepackage {tgschola}
%%: OPTIONAL:  \usepackage {tgtermes}
%%: OPTIONAL:  \usepackage {times}
%%: OPTIONAL:  \usepackage {utopia}

\usepackage {tgbonum}                   % for document font family

%%% ====================================================================
%%% The remaining packages are required by this particular thesis,
%%% but other theses will almost certainly need different packages:
%%%
%%% WARNING: MANY \LaTeX{} packages change dimensions, glue, and/or
%%% formatting styles, and such changes are likely to conflict with
%%% University of Utah Thesis Office requirements.  Therefore, minimize
%%% the number of packages that you include!

\usepackage {amsmath}                   % for AMS math environments
\usepackage {amssymb}                   % for AMS math symbol definitions
%%: OPTIONAL:  \usepackage {amsthm}     % for AMS theorem environments
%%: OPTIONAL:  \usepackage {bm}         % for bold symbol support
\usepackage {bibnames}                  % for BibTeX- and TeX-related names
\usepackage {cite}                      % for numbered and ranged numeric citations
\usepackage {graphicx}                  % for graphics file inclusions
\usepackage {graphpap}                  % for figures in UofUtah sample theses
%%: OPTIONAL:  \usepackage {listings}   % for computer source code listings
%%: OPTIONAL:  \usepackage {longtable}  % for multipage tables
\usepackage {microtype}                 % for improved line and page breaking
%%: OPTIONAL:  \usepackage {multirow}   % for multirow-spanning table cells
%%: OPTIONAL:  \usepackage {pdfpages}   % for inclusion of multipage PDFs
%%: OPTIONAL:  \usepackage {tikz}       % for tikz graphics files
%%: OPTIONAL:  \usetikzlibrary {decorations.markings}
\usepackage {url}                       % for better typesetting of URLs and pathnames
\usepackage {varioref}                  % for improved document cross-references

%%% ====================================================================
%%% Support for showing the sample thesis input and typesetting (NOT
%%% needed by ordinary users)

\usepackage {fancyvrb}                  % for displaying this book's LaTeX input
%%: OPTIONAL:  \usepackage {color}      % for colorizing the sample input
%%: OPTIONAL:  \usepackage {colortbl}                  % for color shading in tables

%%% This is an OPTIONAL package that sets chapter and sectional headings
%%% in color:
%%% Use one or the other of these:
% \usepackage {color}
%%: OPTIONAL:  \usepackage {uuthesis-color}
%%: OPTIONAL:  \definecolor{utahheadingcolor} {rgb}  {0.7, 0.0, 0.0}
%%: OPTIONAL:  \definecolor{utahtheoremcolor} {rgb}  {0.490,0.149,0.804} % purple4
%%: OPTIONAL:  \definecolor{utahtheoremcolor} {rgb}  {0.545,0.137,0.137} % brown4

%%% Here is another, and more convenient, way to define colors, via
%%% aliases of named colors in the X11-derived rgb.sty file

%%: OPTIONAL:  \usepackage{coloralias}
%%: OPTIONAL:  \definecoloralias{utahheadingcolor}{steelblue4}
%%: OPTIONAL:  \definecoloralias{utahtheoremcolor}{hotpink3}

%%% The default heading color is utahred (defined by University Printing
%%% Services as 0.8 red, 0.0 green, 0.0 blue), but you could redefine
%%% that to, for example, a dark blue color, like this AFTER including
%%% the package:
%%%
%%%     \definecolor{utahheading}{rgb}{0,0,0.8}
%%%
%%% NB: Be careful with use of colors in typesetting, and in figures,
%%% because about 6 percent of the human male population is red/green
%%% color blind: they see those colors as shades of brown.  Red and
%%% blue, or blue and green, are better choices for choosing
%%% distinguishable colors.  Also, avoid light colors, especially
%%% yellow, because they are hard to see against white paper and
%%% screen backgrounds, and when printed on black-and-white printers,
%%% where they are rendered in gray, they may be too faint to read.

%%% ====================================================================
%%% Support for a subject index:

%%: OPTIONAL:  \usepackage {uuthesis-index}

%%% ====================================================================
%%% Include latest updates for 2018 version 2.00 base release

\usepackage {uuthesis-updates}          % MANDATORY and LAST of all uuthesis-* packages

%%% ====================================================================
%%% Support for hypertext links in PDF output (this must always be the 
%%% LAST LaTeX package loaded, because it redefines many macros):

%%: OPTIONAL:  \usepackage [pdftex]{hyperref}          % for hypertext links

%%% ====================================================================
%%% This single user-specific file is where all personal customizations
%%% and macro definitions should be placed, and it should come LAST,
%%% after ALL OTHER packages, in case it needs to override some of their
%%% definitions.

\usepackage {mythesis}                  % for personally defined macros

%%% ====================================================================
%%% The student-specific front matter fields are defined here:

\author                 {Kim Buck}
\title                  {Integral Transforms in Quantum Physics}
\thesistype             {thesis}

\dedication             {For my parents, Carol and Dennis.}

%%% Most students need just a short degree name, like this:
\degree                 {Bachelor of Science}
%%% However, multiline degrees are possible, and are done like this:
%%% \degree                 {Doctor of Philosophy \\
%%%                         in \\
%%%                         Mathematical Physics}

%%% For a Master's thesis, use something like this:
%%%
%%% \thesistype             {thesis}
%%%
%%% \degree                 {Master of Science}
%%%
%%% \degree                 {Master of Science \\
%%%                          in \\
%%%                          Experimental Mathematics}
%%%
%%% and omit definitions of \thirdreader and \fourthreader.

%%% College- and Department-level definitions:
\approvaldepartment     {Mathematics}
\department             {Department of Mathematics}
\graduatedean           {Alice B. Toklas}
\departmentchair        {Petrus Marcus Aurelius Featherstone-Hough}

%%% The graduate student's committee members:
\committeechair         {Cornelius L{\'a}nczos}
\firstreader            {Hans Bethe}
%%: \secondreader           {Niels Bohr}
%%: \thirdreader            {Max Born}
%%: \fourthreader           {Paul A. M. Dirac}
\chairtitle             {Professor}

%%% NB: It is rare, but possible, for there to be two chairs, For
%%% example, one student had
%%%
%%% \committeechair{\mbox{\small Andrej Cherkaev and Andrejs Treibergs}}
%%%
%%% The \mbox{} ensures that line breaks cannot happen, and the \small
%%% is necessary to make the names fit on the Dissertation Approval form

%%% Dates that must be adjusted for each academic term, and be permitted
%%% according to the University of Utah Thesis Office:
\submitdate             {May 2020}
\copyrightyear          {2020}

%%% Dates on which committee members approved the thesis
\chairdateapproved      {20 January 2020}
\firstdateapproved      {21 January 2020}
\seconddateapproved     {22 January 2020}
\thirddateapproved      {23 January 2020}
\fourthdateapproved     {24 January 2020}

\usepackage{uuthesis-honors}
\honorsdean             {Sylvia D. Torti, PhD}
\honorsfacultyadvisor   {Sigr{\'\i}{$\eth$}ur Finnbogad{\'o}ttir}

%%% ====================================================================
%%% Typesetting begins here:

\begin{document}

%% Comment out items by inserting a percent % character
\titlepage                              % MANDATORY
%\copyrightpage                          % optional, but strongly recommended
%\dissertationapproval                   % MANDATORY
\preface {abstract} {Abstract}          % optional, but strongly recommended
%%% -*-LaTeX-*-

\chapter* {\LaTeX{} Input for this Thesis}

After stripping comments and empty lines, here is the complete input
\LaTeX{} file that produces this typeset document:

\begin{footnotesize}
\VerbatimInput[numbers = left]{\jobname.stripped}
\end{footnotesize}
                  % ONLY for this sample document
%%% -*-LaTeX-*-

\chapter* {\LaTeX{} Processing for this Thesis}

With the help of definitions and rules in the \texttt{Makefile}, this
thesis is typeset in under \emph{four seconds} with a \emph{single}
operating-system command:

\begin{Verbatim}
    make
\end{Verbatim}

\noindent
That asks the \texttt{make} utility to run these commands to produce a
consistent typeset thesis in Portable Document Format (PDF):

\begin{Verbatim}
    egrep -v '^%|^ *$' honors-thesis-0.tex > honors-thesis-0.stripped

    pdflatex honors-thesis-0.tex
    bibtex honors-thesis-0

    [same command block, repeated 3 more times]
\end{Verbatim}

\noindent
The \texttt{egrep} filtering step is needed to remove empty lines and
comments, producing the bare \LaTeX{} commands shown earlier.  Cross
references in the thesis text, and also inside bibliography entries,
require multiple runs of \texttt{pdflatex} and \texttt{bibtex}.
             % ONLY for this sample document
\dedicationpage                         % optional
\tableofcontents                        % MANDATORY
\settowidth {\fignumwidth} {99.99}
\listoffigures                          % MANDATORY (if there are figures)
\listoftables                           % MANDATORY (if there are tables)
%
%%% Optional front matter page(s), made from source "notation.tex". If
%%% you don't need it, then comment out the \optionalfront command
%%% line!  The first argument is the (unnumbered) section header for
%%% the text supplied by the file input by the second argument; that
%%% file must NOT contain \chapter, \section, \subsection, \ldots{}
%%% sectioning commands.

%%: OPTIONAL:  \optionalfront {Notation and Symbols} {\input{notation}}

%%% Uncomment this if you want the contents of acknowledge.tex typeset
%%% here.  Note that both "Acknowledgement" and "Acknowledgment" are
%%% accepted spellings of that word.

% \preface {acknowledge} {Acknowledgements}

%%% Demonstrations of thesis typesetting features for the sample thesis.
%%% Once you have seen the examples, you can comment out this line.
%%: OPTIONAL: \optionalfront {Typesetting Experiments} {\input{samples.tex}}

\uumaintext     % Start normal page numbering: parts and chapters follow.

%%% -*-LaTeX-*-

%% NB: The command \chapter[short]{long} uses the bracketed short
%% form for two purposes: (a) table of contents, and (b) running
%% page headers.  However, the UofUtah Thesis Office requires the
%% long form in the table of contents, so we must supply the running
%% header text separately in the \chaptermark{} command.
%%
%% It is best to avoid long titles: six or fewer words should suffice.

\chapter[The First,
        Having a Short Alias
        for the Running Page Headers,
        and a Rather Long Definition
        that Takes Multiple Lines with
        User-Defined Line Breaks]
        {The First,                    \\
        Having a Short alias           \\
        for the Running Page Headers,  \\
        and a Rather Long Definition   \\
        that Takes Multiple Lines with \\
        User-Defined Line Breaks}
\chaptermark{The First}

This is a chapter.  Remember  that there should \emph{always}
be at least a few lines of prose after each sectional heading:
failure to do so is a disservice to your readers, and also
produces incorrect vertical spacing.

\section{The first section}

\blah

In \figref{fig1}, we have a picture, and the \LaTeX{} markup
to include it looks like this:
%
\begin{verbatim}
\begin{figure}[t]
    \centerline{\includegraphics{fig1}}
    \caption{The first figure.}%
    \figlabel{fig1}
\end{figure}
\end{verbatim}
%
We intentionally omitted an extension on the filename, so that this
document can be processed with \verb=latex= to get an output
\verb=.dvi= file, or with \verb=pdflatex= to get an output \verb=.pdf=
file.  The first case uses the file \verb=fig1.eps=, and the second
uses \verb=fig1.pdf=.  The \verb=distill= or \verb=ps2pdf= commands
can be used to convert from \emph{Encapulated PostScript}%
files to \emph{Portable Document Format}%
files.

\begin{figure}[t]
    \centerline{\includegraphics{fig1}}
    \caption{The first figure.}%
    \figlabel{fig1}
\end{figure}

\blah

\subsection{The first subsection}

\blah

\subsection{The second subsection}

\blah

\subsection{The third subsection}

\blah

\subsubsection{The first subsubsection}

\blah

\subsubsection{The second subsubsection}

\blah

\paragraph{The first numbered paragraph}

\blah

\paragraph{The second numbered paragraph}

\blah

\section{The second section}

\blah

\blah

\blah

\blah

In \shortfigref{fig2}, we have another picture.

\blah

\begin{figure}[b]
    \centerline{\includegraphics{fig2}}
    \caption{The second figure.}%
    \figlabel{fig2}
\end{figure}

\blah

In \shorttabref{lc-greek}, we show the 24-character lowercase Greek alphabet.

\blah

\blah

\begin{table}[b]
    \caption{Lowercase Greek letters.}%
    \tablabel{lc-greek}
    \vspace{\baselineskip}
    \begin{center}
        \newcommand {\omicron}  {o}
        \begin{tabular}{l@{\hspace{5cm}}l}
            \hline
            $\alpha$                    & alpha \\
            $\beta$                     & beta \\
            $\gamma$                    & gamma \\
            $\delta$                    & delta \\
            $\epsilon$, $\varepsilon$   & epsilon \\
            $\zeta$                     & zeta \\
            $\eta$                      & eta \\
            $\theta$, $\vartheta$       & theta \\
            $\iota$                     & iota \\
            $\kappa$                    & kappa \\
            $\lambda$                   & lambda \\
            $\mu$                       & mu \\
            $\nu$                       & nu \\
            $\xi$                       & xi \\
            $\omicron$                  & omicron \\
            $\pi$                       & pi \\
            $\rho$                      & rho \\
            $\sigma$, $\varsigma$       & sigma \\
            $\tau$                      & tau \\
            $\upsilon$                  & upsilon \\
            $\phi$, $\varphi$           & phi \\
            $\chi$                      & chi \\
            $\psi$                      & psi \\
            $\omega$                    & omega \\
            \hline
        \end{tabular}
    \end{center}
\end{table}

\blah

\blah

\blah

\blah

\blah

\section{The third section}

\blah

In \tabref{uc-greek}, we show the 24-character uppercase Greek
alphabet, 13 of which are identical with Latin letters, because the
Romans borrowed several letters from the earlier Greek alphabet.
However, the letter sounds do not always carry over: notice in
particular the different names of the letter shapes \textbf{H} and
\textbf{P}.  In Modern Greek, {\boldmath $\beta$} is pronounced
\emph{veeta}; the letter pair {\boldmath $\mu\,\pi$} is used to get a
\emph{bee} sound;

\blah

\begin{table}[b]
    \caption[Uppercase Greek letters.]
            {
                Uppercase Greek letters.  Notice that several have the
                same letter shapes as Latin letters, and for those,
                \TeX{} does not define macro names.  For convenience,
                we supply our own definitions of these macros:
                \texttt{\char 92\relax Alpha},
                \texttt{\char 92\relax Beta},
                \texttt{\char 92\relax Epsilon},
                \texttt{\char 92\relax Zeta},
                \texttt{\char 92\relax Eta},
                \texttt{\char 92\relax Iota},
                \texttt{\char 92\relax Kappa},
                \texttt{\char 92\relax Mu},
                \texttt{\char 92\relax Nu},
                \texttt{\char 92\relax Omicron},
                \texttt{\char 92\relax Rho},
                \texttt{\char 92\relax Tau}, and
                \texttt{\char 92\relax Chi}.
            }%
    \tablabel{uc-greek}
    \vspace{\baselineskip}
    \begin{center}
        \newcommand {\Alpha}    {A}
        \newcommand {\Beta}     {B}
        \newcommand {\Epsilon}  {E}
        \newcommand {\Zeta}     {Z}
        \newcommand {\Eta}      {H}
        \newcommand {\Iota}     {I}
        \newcommand {\Kappa}    {K}
        \newcommand {\Mu}       {M}
        \newcommand {\Nu}       {N}
        \newcommand {\Omicron}  {O}
        \newcommand {\Rho}      {P}
        \newcommand {\Tau}      {T}
        \newcommand {\Chi}      {X}
        \begin{tabular}{l@{\hspace{5cm}}l}
            \hline
            $\Alpha$                    & Alpha \\
            $\Beta$                     & Beta \\
            $\Gamma$                    & Gamma \\
            $\Delta$                    & Delta \\
            $\Epsilon$                  & Epsilon \\
            $\Zeta$                     & Zeta \\
            $\Eta$                      & Eta \\
            $\Theta$                    & Theta \\
            $\Iota$                     & Iota \\
            $\Kappa$                    & Kappa \\
            $\Lambda$                   & Lambda \\
            $\Mu$                       & Mu \\
            $\Nu$                       & Nu \\
            $\Xi$                       & Xi \\
            $\Omicron$                  & Omicron \\
            $\Pi$                       & Pi \\
            $\Rho$                      & Rho \\
            $\Sigma$                    & Sigma \\
            $\Tau$                      & Tau \\
            $\Upsilon$                  & Upsilon \\
            $\Phi$                      & Phi \\
            $\Chi$                      & Chi \\
            $\Psi$                      & Psi \\
            $\Omega$                    & Omega \\
            \hline
        \end{tabular}
    \end{center}
\end{table}

\blah

\blah

\blah

\section{Free software packages}

%%: \newcommand {\F} [1] {\texttt{\color{fsfcolor}#1}\index{#1@\protect \fsfname{#1}}}

\newcommand {\F} [1] {\texttt{#1}}

The Free Software Foundation offers almost 300 software packages, most
easily portable to many different operating systems and CPU platforms.
They include at least these:

\begin{sloppypar}
    \F{a2ps},
    \F{acct},
    \F{acm},
    \F{adns},
    \F{alive},
    \F{anubis},
    \F{apl},
    \F{archimedes},
    \F{aris},
    \F{aspell},
    \F{auctex},
    \F{autoconf-archive},
    \F{autoconf},
    \F{autogen},
    \F{automake},
    \F{avl},
    \F{ballandpaddle},
    \F{barcode},
    \F{bash},
    \F{bayonne},
    \F{bc},
    \F{binutils},
    \F{bison},
    \F{bool},
    \F{bpel2owfn},
    \F{c-graph},
    \F{ccaudio},
    \F{ccd2cue},
    \F{ccrtp},
    \F{ccscript},
    \F{cfengine},
    \F{cflow},
    \F{cgicc},
    \F{chess},
    \F{cim},
    \F{classpath},
    \F{classpathx},
    \F{clisp},
    \F{combine},
    \F{commoncpp},
    \F{complexity},
    \F{config},
    \F{coreutils},
    \F{cpio},
    \F{cppi},
    \F{cssc},
    \F{cursynth},
    \F{dap},
    \F{datamash},
    \F{ddd},
    \F{ddrescue},
    \F{dejagnu},
    \F{denemo},
    \F{dico},
    \F{diction},
    \F{diffutils},
    \F{dionysus},
    \F{direvent},
    \F{dismal},
    \F{dominion},
    \F{easejs},
    \F{ed},
    \F{edma},
    \F{electric},
    \F{emacs},
    \F{emms},
    \F{enscript},
    \F{fdisk},
    \F{ferret},
    \F{findutils},
    \F{fisicalab},
    \F{flex},
    \F{fontutils},
    \F{freedink},
    \F{freefont},
    \F{freeipmi},
    \F{gama},
    \F{garpd},
    \F{gawk},
    \F{gcal},
    \F{gcc},
    \F{gcide},
    \F{gcl},
    \F{gcompris},
    \F{gdb},
    \F{gdbm},
    \F{gengen},
    \F{gengetopt},
    \F{gettext},
    \F{gforth},
    \F{ggradebook},
    \F{ghostscript},
    \F{gift},
    \F{gleem},
    \F{glibc},
    \F{global},
    \F{glpk},
    \F{gmp},
    \F{gnash},
    \F{gnats},
    \F{gnatsweb},
    \F{gnu-c-manual},
    \F{gnu-crypto},
    \F{gnu-pw-mgr},
    \F{gnubatch},
    \F{gnubik},
    \F{gnucap},
    \F{gnucobol},
    \F{gnudos},
    \F{gnue},
    \F{gnugo},
    \F{gnuit},
    \F{gnujump},
    \F{gnukart},
    \F{gnumach},
    \F{gnun},
    \F{gnunet},
    \F{gnupod},
    \F{gnuprologjava},
    \F{gnuradio},
    \F{gnurobots},
    \F{gnuschool},
    \F{gnushogi},
    \F{gnusound},
    \F{gnuspeech},
    \F{gnuspool},
    \F{gnustep},
    \F{gnutls},
    \F{gnutrition},
    \F{gnuzilla},
    \F{goptical},
    \F{gperf},
    \F{gprolog},
    \F{greg},
    \F{grep},
    \F{groff},
    \F{grub},
    \F{gsasl},
    \F{gsegrafix},
    \F{gsl},
    \F{gslip},
    \F{gsrc},
    \F{gss},
    \F{gtypist},
    \F{guile-gnome},
    \F{guile-gtk},
    \F{guile-ncurses},
    \F{guile-opengl},
    \F{guile-rpc},
    \F{guile-sdl},
    \F{guile},
    \F{gv},
    \F{gvpe},
    \F{gxmessage},
    \F{gzip},
    \F{halifax},
    \F{health},
    \F{hello},
    \F{help2man},
    \F{hp2xx},
    \F{httptunnel},
    \F{hurd},
    \F{hyperbole},
    \F{idutils},
    \F{ignuit},
    \F{indent},
    \F{inetutils},
    \F{intlfonts},
    \F{jacal},
    \F{jel},
    \F{jwhois},
    \F{kawa},
    \F{less},
    \F{libcdio},
    \F{libextractor},
    \F{libffcall},
    \F{libiconv},
    \F{libidn},
    \F{libmatheval},
    \F{libmicrohttpd},
    \F{librejs},
    \F{libsigsegv},
    \F{libtasn1},
    \F{libtool},
    \F{libunistring},
    \F{libxmi},
    \F{lightning},
    \F{lilypond},
    \F{liquidwar6},
    \F{lsh},
    \F{m4},
    \F{macchanger},
    \F{mailman},
    \F{mailutils},
    \F{make},
    \F{marst},
    \F{maverik},
    \F{mc},
    \F{mcron},
    \F{mcsim},
    \F{mdk},
    \F{metahtml},
    \F{mifluz},
    \F{mig},
    \F{miscfiles},
    \F{mit-scheme},
    \F{moe},
    \F{motti},
    \F{mpc},
    \F{mpfr},
    \F{mpria},
    \F{mtools},
    \F{myserver},
    \F{nano},
    \F{ncurses},
    \F{nettle},
    \F{non-gnu},
    \F{ocrad},
    \F{octave},
    \F{oleo},
    \F{orgadoc},
    \F{osip},
    \F{paperclips},
    \F{parallel},
    \F{parted},
    \F{patch},
    \F{pem},
    \F{pexec},
    \F{phantom},
    \F{pies},
    \F{plotutils},
    \F{proxyknife},
    \F{pspp},
    \F{psychosynth},
    \F{pth},
    \F{pyconfigure},
    \F{radius},
    \F{rcs},
    \F{readline},
    \F{recutils},
    \F{reftex},
    \F{remotecontrol},
    \F{rottlog},
    \F{rpge},
    \F{rush},
    \F{sather},
    \F{sauce},
    \F{savannah},
    \F{scm},
    \F{screen},
    \F{sed},
    \F{serveez},
    \F{sharutils},
    \F{shishi},
    \F{shmm},
    \F{shtool},
    \F{sipwitch},
    \F{slib},
    \F{smalltalk},
    \F{solfege},
    \F{spacechart},
    \F{spell},
    \F{sqltutor},
    \F{src-highlite},
    \F{stow},
    \F{superopt},
    \F{swbis},
    \F{tar},
    \F{termcap},
    \F{termutils},
    \F{teseq},
    \F{teximpatient},
    \F{texinfo},
    \F{thales},
    \F{time},
    \F{tramp},
    \F{trueprint},
    \F{unifont},
    \F{units},
    \F{unrtf},
    \F{userv},
    \F{uucp},
    \F{vc-dwim},
    \F{vcdimager},
    \F{vera},
    \F{wb},
    \F{wdiff},
    \F{websocket4j},
    \F{wget},
    \F{which},
    \F{windows},
    \F{xaos},
    \F{xboard},
    \F{xhippo},
    \F{xlogmaster},
    \F{xnee},
    \F{xorriso}, and
    \F{zile}.
\end{sloppypar}

\blah

In \shortfigref{fig3}, we have yet another picture.

\blah

\blah

\blah

\blah

\begin{figure}[b]
    \centerline{\includegraphics{fig3}}
    \caption[The third figure.]
            {%
                The third figure.  This one has both short and long captions.
                \blah \par
                \blah \par
                \blah
            }%
    \figlabel{fig3}
\end{figure}

\blah

\blah

\blah

\blah

\blah

\blah

\blah

\blah


\section{Resizing figures}

In \shortfigref{fig4} through \figref{fig8}, we show how graphics files
can be rescaled to convenient sizes, with input like this:

\begin{figure}[p]
    \centerline{\includegraphics[scale = 0.5]{fig1}}
    \caption{The fourth figure (at 50\% scale).}%
    \figlabel{fig4}
\end{figure}

\begin{figure}[p]
    \centerline{\includegraphics[scale = 0.75]{fig1}}
    \caption{The fifth figure (at 75\% scale).}%
    \figlabel{fig5}
\end{figure}

\begin{figure}[p]
    \centerline{\includegraphics{fig1}}
    \caption{The sixth figure (at native size).}%
    \figlabel{fig6}
\end{figure}

\begin{figure}[p]
    \centerline{\includegraphics[scale = 1.25]{fig1}}
    \caption{The seventh figure (at 125\% scale).}%
    \figlabel{fig7}
\end{figure}

\begin{figure}[p]
    \centerline{\includegraphics[scale = 1.75]{fig1}}
    \caption{The eighth figure (at 175\% scale).}%
    \figlabel{fig8}
\end{figure}

\begin{verbatim}
\begin{figure}[p]
    \centerline{\includegraphics[scale = 0.5]{fig1}}
    \caption{The fourth figure (at 50\% scale).}%
    \figlabel{fig4}
\end{figure}

\begin{figure}[p]
    \centerline{\includegraphics[scale = 0.75]{fig1}}
    \caption{The fifth figure (at 75\% scale).}%
    \figlabel{fig5}
\end{figure}

\begin{figure}[p]
    \centerline{\includegraphics{fig1}}
    \caption{The sixth figure (at native size).}%
    \figlabel{fig6}
\end{figure}

\begin{figure}[p]
    \centerline{\includegraphics[scale = 1.25]{fig1}}
    \caption{The seventh figure (at 125\% scale).}%
    \figlabel{fig7}
\end{figure}

\begin{figure}[p]
    \centerline{\includegraphics[scale = 1.75]{fig1}}
    \caption{The eighth figure (at 175\% scale).}%
    \figlabel{fig8}
\end{figure}
\end{verbatim}

You can include multiple images, each with its own caption inside a
single \emph{unbreakable} \texttt{figure} environment, like this
example shown in \shortfigref{fig9} and \figref{fig10}, although
you might want to adjust interfigure vertical space with a
\verb=\vspace{}= command:

\begin{verbatim}
\begin{figure}[t]
    \centerline{\includegraphics[scale = 0.5]{fig1}}
    \caption{The fourth figure (at 50\% scale).}%
    \figlabel{fig9}
    \vspace{3ex}
    \centerline{\includegraphics[scale = 0.75]{fig1}}
    \caption{The fifth figure (at 75\% scale).}%
    \figlabel{fig10}
\end{figure}
\end{verbatim}

\begin{figure}[t]
    \centerline{\includegraphics[scale = 0.5]{fig1}}
    \caption[The ninth figure (at 50\% scale)]%
            {The ninth figure (at 50\% scale), boxed with the tenth figure.}
    \figlabel{fig9}
    \vspace{3ex}
    \centerline{\includegraphics[scale = 0.75]{fig1}}
    \caption[The tenth figure (at 75\% scale)]%
            {The tenth figure (at 75\% scale), boxed with the ninth figure.}%
    \figlabel{fig10}
\end{figure}

\blah

\blah

\blah

As a final example in this chapter, \figref{picture-mode}
shows how you can use \LaTeX{} \texttt{picture} mode for
annotating and positioning graphics images prepared outside
\LaTeX{}.  The input that produced that figure looks like
this:

%
\begin{verbatim}
\begin{figure}[t]
    %% The original image is 216bp wide by 72bp high, but we
    %% rescale it to 150 picture units divided by \unitlength:
    %% 150 / 0.75 = 112.5 mm
    \newcommand {\myfig} {\includegraphics[width = 112.5mm]{fig1}}

    \begin{center}
        %% The \unitlength is chosen to make the complete picture fit
        %% within the page margins

        \setlength{\unitlength}{0.75mm}

        %%%     insert (width,height)(lower-left-x,lower-left-y)
        \begin{picture}(170,70)(10,10)
            %% Place the included image FIRST!
            \put(10,10) {\myfig}

            %% Everything that follows OVERLAYS the original image!

            \graphpaper[10](0,0)(170,70)

            %% Mark the image center and corners by centered bullets
            \newcommand {\thedot} {\makebox (0,0) {$\bullet$}}
            \put( 85, 35) {\thedot}
            \put( 10, 10) {\thedot}
            \put( 10, 60) {\thedot}
            \put(160, 10) {\thedot}
            \put(160, 60) {\thedot}

            \put( 10, 10) {\makebox (0,0) [r] {lower-left}}
            \put(160, 10) {\makebox (0,0) [l] {lower-right}}
            \put( 10, 60) {\makebox (0,0) [r] {upper-left}}
            \put(160, 60) {\makebox (0,0) [l] {upper-right}}
        \end{picture}
    \end{center}

    \vspace{2\baselineskip}

    \caption[Using \LaTeX{} \texttt{picture} mode]
            {Using \LaTeX{} \texttt{picture} mode for figure labeling
             and positioning.}
    \figlabel{picture-mode}
\end{figure}
\end{verbatim}

\begin{figure}[t]
    %% The original image is 216bp wide by 72bp high, but we
    %% rescale it to 150 picture units divided by \unitlength:
    %% 150 / 0.75 = 112.5 mm
    \newcommand {\myfig} {\includegraphics[width = 112.5mm]{fig1}}

    \begin{center}
        %% The \unitlength is chosen to make the complete picture fit
        %% within the page margins

        \setlength{\unitlength}{0.75mm}

        %%%     insert (width,height)(lower-left-x,lower-left-y)
        \begin{picture}(170,70)(10,10)
            %% Place the included image FIRST!
            \put(10,10) {\myfig}

            %% Everything that follows OVERLAYS the original image!

            \graphpaper[10](0,0)(170,70)

            %% Mark the image center and corners by centered bullets
            \newcommand {\thedot} {\makebox (0,0) {$\bullet$}}
            \put( 85, 35) {\thedot}
            \put( 10, 10) {\thedot}
            \put( 10, 60) {\thedot}
            \put(160, 10) {\thedot}
            \put(160, 60) {\thedot}

            \put( 10, 10) {\makebox (0,0) [r] {lower-left}}
            \put(160, 10) {\makebox (0,0) [l] {lower-right}}
            \put( 10, 60) {\makebox (0,0) [r] {upper-left}}
            \put(160, 60) {\makebox (0,0) [l] {upper-right}}
        \end{picture}
    \end{center}

    \vspace{2\baselineskip}

    \caption[Using \LaTeX{} \texttt{picture} mode]
            {Using \LaTeX{} \texttt{picture} mode for figure labeling
             and positioning.}
    \figlabel{picture-mode}
\end{figure}

%
\section{Summary and conclusions}

\blah

\blah
\blah

\blah

%%% Index phrases should be attached to an important word of a phrase,
%%% and are usually best kept on a separate line by terminating the
%%% previous line with a percent comment without intervening space, as
%%% in this example:
%%%
%%%     \newcommand {\X} [1] {#1\index{#1}}
%%%
%%%     African ungulates,%
%%%     \index{African ungulate}
%%%     like the \X{gnu}, \X{impala}, \X{kudu}, and \X{springbok}
%%%     live mostly in hot climate and consume vegetation.
%%%
%%% However, for this document, we only want lots of index entries to
%%% populate a sample topic index.

%%% ====================================================================
%%% Cross-references for index entries should be specified only once:

%%% -*-LaTeX-*-

\chapter{The Second}

This is a chapter.

\blah

%%% -*-LaTeX-*-

\chapter{The Third}

This is a chapter.

\blah

%%% -*-LaTeX-*-

\chapter{The Fourth}

This is a chapter.

\blah

\blah

\blah

\section{More on the topic}

\blah

\blah

\section{Even more on the topic}

\blah

\blah

\section{Summary and conclusions}

\blah


\appendix       % Chapters, sections are now appendix style A, A.1, A.2, B, C, D, ...

%%% -*-LaTeX-*-

\chapter{The First}

This is an appendix.  The \LaTeX{} markup for an appendix
is, surprisingly, \verb=\chapter=.  The \verb=\appendix= command does
not produce a heading; instead, it just changes the numbering style
from numeric to alphabetic, and it changes the heading prefix from
\textbf{Chapter} to \textbf{Appendix}.

\blah \blah \blah

%%% -*-LaTeX-*-

\chapter{The Second}

This is an appendix.

\blah \blah \blah

%%% -*-LaTeX-*-

\chapter{The Third}

This is an appendix.

There are several books
\cite{%
    Singh:1997:FEE,%
    Salomon:1995:AT,%
    Robbins:2005:CSS,%
    Randell:1982:ODC,%
    Olver:2010:NHM,%
    Mittelbach:2004:LC,%
    Lamport:1985:LDP,%
    Knuth:1999:DT,%
    Knuth:1986:TB,%
    Knuth:1986:MB,%
    Farmelo:2009:SMH%
}
listed in our bibliography.

We also reference several journal articles
\cite{%
    Wiles:1995:MEC,%
    Taylor:1995:RTP,%
    Lahiri:2002:ADA,%
    Kim:1999:AFC,%
    Johnson:1978:LDT,%
    Huskey:1980:LLC,%
    Heilbron:1969:GBA,%
    Hall:1994:PNE,%
    Goldstine:1946:ENI,%
    Einstein:1911:BMAb,%
    Einstein:1911:BMAa,%
    Einstein:1906:NBM,%
    Cody:1981:APF,%
    Beebe:1989:PCP,%
    Babbage:1910:BBA%
}
and three famous doctoral theses of later winners
\cite{%
    Einstein:1905:NBM,%
    Dirac:1926:QM,%
    Bohr:1911:SME%
}
of the Nobel Prize in Physics (1922, 1933, and
1921):

Notice that, even though those citations appeared
in \LaTeX{} \verb=\cite{...}= commands with their
\BibTeX{} citation labels in reverse alphabetical
order, thanks to the \verb=cite= package,
their reference-list numbers have been sorted in
numerically ascending order, and then
range-reduced.

Mention should also be made of a famous Dutch
computer scientist's first publication
\cite{Dijkstra:1953:FBV}.

Font metrics are an important, albeit low-level,
aspect of typesetting. See the \emph{Adobe
Systems} manual about that company's procedures
\cite{Adobe:1990:AFM}.

The bibliography at the end of this thesis contains several examples
of documents with non-English titles, and their \BibTeX{} entries
provide title translations following the practice recommended by the
American Mathematical Society and SIAM\@.  Here is a sample entry that
shows how to do so:
%
\begin{verbatim}
@PhdThesis{Einstein:1905:NBM,
  author =       "Albert Einstein",
  title =        "{Eine Neue Bestimmung der Molek{\"u}ldimensionen}.
                 ({German}) [{A} new determination of molecular
                 dimensions]",
  type =         "Inaugural dissertation",
  school =       "Bern Wyss.",
  address =      "Bern, Switzerland",
  year =         "1905",
  bibdate =      "Fri Dec 17 10:46:57 2004",
  bibsource =    "http://www.math.utah.edu/pub/tex/bib/einstein.bib",
  note =         "Published in \cite{Einstein:1906:NBM}.",
  acknowledgement = ack-nhfb,
  language =     "German",
  advisor =      "Alfred Kleiner (24 April 1849--3 July 1916)",
  URL =          "http://en.wikipedia.org/wiki/Alfred_Kleiner",
  remark =       "Received August 19, 1905 and published February 8,
                 1906.",
  Schilpp-number = "6",
}
\end{verbatim}

The \texttt{note} field in that entry refers to another bibliography
entry that need not have been directly cited in the document text.
Such cross-references are common in \BibTeX{} files, especially for
journal articles where there may be later comments and corrigenda that
should be mentioned.  Embedded \verb=\cite{}= commands ensure that
those possibly important other entries are always included in the
reference list when the entry is cited.  The last bibliography entry
\cite{Wiles:1995:MEC} in this thesis has a long \texttt{note} field
that tells more about what some may view as the most important paper
in mathematics in the last century.

When entries cite other entries that cite other entries that cite
other entries that \ldots{}, multiple passes of \LaTeX{} and \BibTeX{}
are needed to ensure consistency.  That is another reason why document
compilation should be guided by a \texttt{Makefile} or a batch script,
rather than expecting the user to remember just how many passes are
needed.

\BibTeX{} entries are \emph{extensible}, in that arbitrary key\slash
value pairs may be present that are not necessarily recognized by any
bibliography style files.  The \texttt{advisor},
\texttt{acknowledgement}, \texttt{bibdate}, \texttt{bibsource},
\texttt{language}, \texttt{remark}, and \texttt{Schilpp-number} fields
are examples, and may be used by other software that processes
\BibTeX{} entries, or by humans who read the entries.  \texttt{DOI}
and \texttt{URL} fields are currently recognized by only a few styles,
but that situation will likely change as publishers demand that such
important information be included in reference lists.

In \BibTeX{}, \texttt{title} and thesis \texttt{type} 
fields, braces protect words, such as
proper nouns and acro\-nyms, that cannot be downcased if the selected
bibliography style would otherwise do so.  In German, all nouns are
capitalized, and the simple way to ensure their protection is to brace
the entire German text in the title, as we did in the entry above.

The world's first significant computer program may
have been that written in 1842 by Lady Augusta Ada
Lovelace (1815--1852) for the computation of Bernoulli
numbers \cite{Huskey:1980:LLC,Kim:1999:AFC}.  She
was the assistant to Charles Babbage
(1791--1871), and they are the world's first
computer programmers. The programming language
\emph{Ada} is named after her, and is defined in
the ANSI/MIL-STD-1815A Standard; its number
commemorates the year of her birth.

We do not discuss mathematical \emph{transforms}
in this dissertation, but you can find that phrase
in the index (except that this sample thesis doesn't have one!)

\blah

\blah
\blah

\blah


%%% The choice of bibliography style is a major decision, jointly made
%%% by you, your thesis advisor and the thesis editor. Common choices are
%%% one of the four standard BibTeX styles (abbrv, alpha, plain, and unsrt),
%%% or enhanced styles like acm, amsplain, siam, and hundreds of others
%%% available in TeX Live, and other Unix and Windows TeX distributions.
%%%
%%% Do NOT handcode your reference list, because you are unlikely to
%%% achieve consistency or conformance to the University of Utah Thesis Office
%%% requirements: let BibTeX do that tedious job for you!
%%%
%%% Remember that reference-list metadata in BibTeX files remains
%%% constant across journals and publishers, and is are often reused
%%% in other documents and shared with others, whereas formatted
%%% reference-list styles change: with BibTeX, you only need to record
%%% the metadata once.
%%%
%%% If you prefer named, rather than numeric or tagged citations, you
%%% may use styles such as authordate{1,2,3,4}, chicago, harvard, or
%%% natbib.  Be aware, however, that most of those require an
%%% additional \usepackage{} command to supply \LaTeX{} with
%%% definitions of commands that the style needs, and that there are
%%% usually several flavors of LaTeX citation commands beyond the
%%% standard \cite{} command that you need to understand before you
%%% can use them properly in your prose.

%%% This tells BibTeX to read siam.bst from the first directory where
%%% it is found in the BSTINPUTS search path:

\bibliographystyle{plain}

%%% This can also specify a comma-separated list (without embedded
%%% spaces) of *.bib files found by BibTeX in its BIBINPUTS search
%%% path.  The argument \jobname means the base name of the top-level
%%% LaTeX file, avoiding an unnecessary filename dependence here.
%%%
%%% BibTeX writes only one .bbl file, no matter how many *.bib files
%%% are listed here, using the name \jobname.bbl.
%%%
%%% LaTeX reads BibTeX's formatted reference list from the file
%%% \jobname.bbl.

\bibliography{\jobname}

\end {document}
